% Minimal preamble for Open Logic Project sources (pinned commit 1e960beff9ed7835bf3e3f1335e21af3439cd107).
%
% OLP's own master preamble cannot run on the project host: TeX Live 2025 there lacks
% prooftree.sty, ntheorem.sty, extarrows.sty, mathdots.sty, isomath.sty and subfiles.sty.
% Following the project's own rule ("if the author preamble needs packages the toolchain lacks,
% do not invent macros -- use a minimal working preamble"), this file declares, verbatim in
% spirit and numbering behaviour, exactly the theorem-like environments OLP defines in
% sty/open-logic-book-envs.sty (thm/lem/prop/cor/defn/ex/prob/rem/note/case/axiom/conv),
% plus the shorthand accents OLP uses (`!!a`, `!!b`, `!!{...}`).
%
% Author text is untouched; only the packaging is ours.
\documentclass[10pt,a4paper]{article}

\usepackage{amsmath}
\usepackage{amssymb}
\usepackage{amsthm}
\usepackage{graphicx}
\usepackage{hyperref}
\usepackage{enumitem}
\usepackage[framemethod=default]{mdframed}

% ---- OLP shorthand accents (from sty/open-logic.sty) ----
% NB: \! is taken by LaTeX (negative thin space), so we only define the !!a / !!b families
% and the \olup helper used by them.
\newcommand*{\olup}[1]{\MakeUppercase{#1}}
\def\!!!a#1{\olup{#1}}
\def\!!a#1{\olup{#1}}
\def\!!b#1{\olup{#1}}

% ---- theorem-like environments, same family numbering as OLP ----
\theoremstyle{plain}
\newtheorem{thm}{Theorem}[section]
\newtheorem{lem}[thm]{Lemma}
\newtheorem{prop}[thm]{Proposition}
\newtheorem{cor}[thm]{Corollary}
\newtheorem{axiom}{Axiom}

\theoremstyle{definition}
\newtheorem{defn}[thm]{Definition}
\newtheorem{ex}[thm]{Example}
\newtheorem{prob}{Problem}[section]
\newtheorem{probd}{Problem}
\newtheorem{conv}[thm]{Convention}

\theoremstyle{remark}
\newtheorem{rem}[thm]{Remark}
\newtheorem{note}[thm]{Note}
\newtheorem{case}[thm]{Case}

% ---- environments OLP defines as wrappers ----
\newenvironment{probtag}[1]{\begin{prob}}{\end{prob}}
\newenvironment{explain}{}{}
\newenvironment{explain*}{}{}
\newenvironment{editorial}{}{}
\newenvironment{digress}{}{}
\newenvironment{digress*}{}{}
\newenvironment{history}{}{}
\newenvironment{reference}{}{}
\newcommand*{\olfileid}[3]{}
\newcommand*{\olsection}[1]{\section*{#1}}
\newcommand*{\olchapter}[3]{}
\newcommand*{\olpart}[2]{}
\newcommand*{\olimport}[1]{}
\newcommand*{\olinclude}[1]{}
\newcommand*{\olpath}[1]{}
\newcommand*{\ollangid}[1]{}
\newcommand*{\problemtag}[1]{}
\newcommand*{\problemauthor}[1]{}
\newcommand*{\ollabel}[1]{\label{#1}}
% \olref takes 0-3 OPTIONAL arguments plus one mandatory label:
%     \olref{lab}  |  \olref[from]{lab}  |  \olref[from][to]{lab}  |  \olref[a][b][c]{lab}
% A single-argument definition silently swallows "[sth]" and emits a broken reference
% (measured on the host: the text came out as "??iso]x"). The chain below consumes the
% brackets and cross-references only the mandatory argument.
\makeatletter
\newcommand{\olref}{\@ifnextchar[{\olref@i}{\olref@final}}
\def\olref@i[#1]{\@ifnextchar[{\olref@ii}{\olref@final}}
\def\olref@ii[#1]{\@ifnextchar[{\olref@iii}{\olref@final}}
\def\olref@iii[#1]{\olref@final}
\def\olref@final#1{\ref{#1}}
\makeatother
\newcommand*{\olnewpage}{\clearpage}

% ===================== OLP compatibility layer (1) =====================

\makeatletter

% ---------------- proof trees (OLP's prooftree is a bussproofs derivative) ----------------
\providecommand{\AxiomC}[1]{}
\providecommand{\UnaryInfC}[1]{}
\providecommand{\BinaryInfC}[1]{}
\providecommand{\TrinaryInfC}[1]{}
\providecommand{\QuaternaryInfC}[1]{}
\providecommand{\QuinaryInfC}[1]{}
\providecommand{\RightLabel}[1]{}
\providecommand{\LeftLabel}[1]{}
\providecommand{\doubleLine}{}
\providecommand{\noLine}{}
\providecommand{\DisplayProof}{}
\providecommand{\fCenter}{}
\newenvironment{prooftree}{}{}
\newenvironment{ltprooftree}{}{}
\newenvironment{oltableau}{}{}

% ---------------- logical notation ----------------
\providecommand{\lif}{\mathbin{\rightarrow}}
\providecommand{\lforall}{\mathord{\forall}}
\providecommand{\lexists}{\mathord{\exists}}
\providecommand{\ltrue}{\top}
\providecommand{\lfalse}{\bot}
\providecommand{\Ltrue}{\top}
\providecommand{\Lfalse}{\bot}
\providecommand{\Sequent}{\mathrel{\Rightarrow}}
\providecommand{\Proves}{\mathrel{\vdash}}
\providecommand{\Entails}{\mathrel{\models}}
\providecommand{\Subst}[3]{#1[#2/#3]}
\providecommand{\sFmla}[2]{#1(#2)}
\providecommand{\pFmla}[2]{#1(#2)}
\providecommand{\formula}{\relax}
\providecommand{\True}{\top}
\providecommand{\False}{\bot}
\providecommand{\Args}[1]{#1}
\providecommand{\Ts}{\ensuremath{\top}}
\providecommand{\Fi}{\ensuremath{\bot}}
\providecommand{\lfalsum}{\bot}

% ---------------- objects, structures, metatheory ----------------
\providecommand{\Obj}{}
\providecommand{\Struct}{\ensuremath{\mathfrak{A}}}
\providecommand{\Domain}{\ensuremath{|\mathfrak{A}|}}
\providecommand{\Assign}{\ensuremath{s}}
\providecommand{\Sat}[2]{\ensuremath{#1 \models #2}}
\providecommand{\mSat}[3]{\ensuremath{#1 \models #2[#3]}}
\providecommand{\mModel}[2]{\ensuremath{#1 \models #2}}
\providecommand{\Lang}{\ensuremath{\mathcal{L}}}
\providecommand{\Th}{\ensuremath{\mathrm{Th}}}
\providecommand{\Log}{\ensuremath{\mathrm{L}}}
\providecommand{\tuple}[1]{\langle #1 \rangle}
\providecommand{\num}[1]{\underline{#1}}
\providecommand{\cfind}[1]{\varphi_{#1}}
\providecommand{\Nat}{\ensuremath{\mathbb{N}}}
\providecommand{\Mult}{\ensuremath{\mathrm{Mult}}}
\providecommand{\Add}{\ensuremath{\mathrm{Add}}}
\providecommand{\Exp}{\ensuremath{\mathrm{Exp}}}
\providecommand{\Zero}{\ensuremath{0}}
\providecommand{\Succ}{\ensuremath{\mathrm{S}}}
\providecommand{\lambd}{\ensuremath{\lambda}}
\providecommand{\bredpar}{\mathrel{\rightarrow_\beta}}
\providecommand{\bcd}[1]{\ensuremath{\ulcorner #1 \urcorner}}
\providecommand{\Set}[1]{\ensuremath{\{#1\}}}
\providecommand{\indcase}[1]{\par\noindent\emph{Case #1}.\quad}
\providecommand{\indfrm}[1]{\par\noindent\emph{Form #1}.\quad}
\providecommand{\iftag}[1]{#1}
\providecommand{\tagitem}[1]{\item[#1]}
\providecommand{\Elim}[1]{}
\providecommand{\Intro}[1]{}
\providecommand{\LeftR}[1]{}
\providecommand{\RightR}[1]{}
\providecommand{\TAss}[1]{}
\providecommand{\TRule}[1]{}

\makeatother
\providecommand{\olpath}{.}

% ===================== OLP compatibility layer (2) =====================
% --- second compatibility block: notation and shorthands used by the frozen OLP units ---
\providecommand{\CutCS}{\mathrm{Cut}}
\providecommand{\LeftSubproofLabel}{\relax}
\providecommand{\RightSubproofLabel}{\relax}
\providecommand{\maeh}{\relax}
\providecommand{\cheight}{\relax}
\providecommand{\ande}{\mathrm{and}_E}
\providecommand{\andi}{\mathrm{and}_I}
\providecommand{\ore}{\mathrm{or}_E}
\providecommand{\ori}{\mathrm{or}_I}
\providecommand{\typeof}{\mathord{:}}
\providecommand{\pValue}[1]{#1}
\providecommand{\Sent}{\mathrm{Sent}}
\providecommand{\cardplus}{\mathbin{+}}
\providecommand{\cardtimes}{\mathbin{\cdot}}
\providecommand{\cardexpo}{\mathbin{\hat{\ }}}
\providecommand{\cardnless}{\not<}
\providecommand{\cardle}{\leq}
\providecommand{\funfromto}[3]{#1\colon #2 \to #3}
\providecommand{\funimage}[1]{\operatorname{im}(#1)}
\providecommand{\fdefined}{\mathord{\downarrow}}
\providecommand{\canonord}{\operatorname{ord}}
\providecommand{\ordtimes}{\mathbin{\cdot}}
\providecommand{\ordeq}[2]{#1 = #2}
\providecommand{\ordneq}[2]{#1 \neq #2}
\providecommand{\ordtype}[1]{\operatorname{ord}(#1)}
\providecommand{\ordsucc}{\mathord{+1}}
\providecommand{\ordunion}{\cup}
\providecommand{\onesphere}{S^1}
\providecommand{\rotationsgroup}{\mathrm{SO}(2)}
\providecommand{\equivrep}[1]{[#1]}
\providecommand{\rlexless}{\prec}
\providecommand{\ZFC}{\mathrm{ZFC}}
\providecommand{\ZF}{\mathrm{ZF}}
\providecommand{\Rat}{\mathbb{Q}}
\providecommand{\Real}{\mathbb{R}}
\providecommand{\PosInt}{\mathbb{Z}^{+}}
\providecommand{\Id}{\mathrm{id}}
\providecommand{\Bin}{\mathrm{Bin}}
\providecommand{\pto}{\rightharpoonup}
\providecommand{\dom}{\operatorname{dom}}
\providecommand{\ran}{\operatorname{ran}}
\providecommand{\len}{\operatorname{len}}
\providecommand{\concat}{\mathbin{{}^\frown}}
\providecommand{\card}[1]{\lvert #1 \rvert}
\providecommand{\Pow}{\mathcal{P}}
\providecommand{\Setabs}[2]{\{\, #1 \mid #2 \,\}}
\providecommand{\comp}{\circ}
\providecommand{\Char}{\chi}
\providecommand{\Prf}{\operatorname{Prf}}
\providecommand{\Prov}{\operatorname{Prov}}
\providecommand{\Proj}{\operatorname{proj}}
\providecommand{\Gn}{\operatorname{gn}}
\providecommand{\gn}{\operatorname{gn}}
\providecommand{\Value}[1]{\operatorname{val}(#1)}
\providecommand{\Deduce}{\mathrm{D}}
\providecommand{\Dual}{\mathrm{d}}
\providecommand{\Ax}{\mathrm{Ax}}
\providecommand{\pSat}[3]{#1 \models #2\,[#3]}
\providecommand{\pAssign}{\sigma}
\providecommand{\bforall}{\forall}
\providecommand{\bexists}{\exists}
\providecommand{\liff}{\mathbin{\leftrightarrow}}
\providecommand{\FV}{\operatorname{FV}}
\providecommand{\UnaryInf}{\ensuremath{\mid}}
\providecommand{\BinaryInf}{\ensuremath{\mid}}
\providecommand{\depth}{\operatorname{depth}}
\providecommand{\oliflabeldef}[3]{#1}
\providecommand{\TMblank}{\sqcup}
\providecommand{\TMstroke}{\mid}
\providecommand{\TMstay}{\mathrm{S}}
\providecommand{\TMright}{\mathrm{R}}
\providecommand{\TMleft}{\mathrm{L}}
\providecommand{\TMendtape}{\triangleright}
\providecommand{\umin}{\mu}
\providecommand{\Sum}{\sum}
\providecommand{\citep}[1]{\cite{#1}}
\providecommand{\citet}[1]{\cite{#1}}
\providecommand{\citeauthor}[1]{\cite{#1}}
\providecommand{\citeyear}[1]{\cite{#1}}

% ===================== OLP compatibility layer (3) =====================
% --- third compatibility block: remaining notation seen in the frozen OLP units ---
\providecommand{\iso}{\cong}
\providecommand{\nszero}{\mathrm{0}}
\providecommand{\nssucc}{\mathrm{S}}
\providecommand{\nsplus}{\mathbin{+}}
\providecommand{\nssless}{<}
\providecommand{\nsless}{<}
\providecommand{\PL}{\mathrm{PL}}
\providecommand{\Nec}{\mathord{\Box}}
\providecommand{\RK}{\mathrm{RK}}
\providecommand{\ST}{\mathrm{ST}}
\providecommand{\SSubst}[3]{#1[#2/#3]}
\providecommand{\subst}[3]{#1[#2/#3]}
\providecommand{\tagfalse}{\mathrm{F}}
\providecommand{\cutr}{\mathrm{Cut}}
\providecommand{\cardneq}{\neq}
\providecommand{\issucc}[1]{\operatorname{succ}(#1)}
\providecommand{\tup}[1]{\langle #1 \rangle}
\providecommand{\Var}{\mathrm{Var}}
\providecommand{\Bew}{\operatorname{Bew}}
\providecommand{\Con}{\operatorname{Con}}
\providecommand{\Cons}{\operatorname{Cons}}
\providecommand{\pDeduce}{\mathrm{D}}
\providecommand{\psat}{\models}
\providecommand{\den}[1]{\llbracket #1 \rrbracket}
\providecommand{\Sem}[2]{\llbracket #1 \rrbracket_{#2}}
\providecommand{\DomainOf}[1]{\lvert #1 \rvert}
\providecommand{\Form}{\mathrm{Form}}
\providecommand{\Terms}{\mathrm{Term}}
\providecommand{\Atom}{\mathrm{Atom}}
\providecommand{\FVars}[1]{\operatorname{FV}(#1)}
\providecommand{\fv}[1]{\operatorname{fv}(#1)}
\providecommand{\tuplep}[1]{\langle #1 \rangle}
\providecommand{\restrict}[2]{#1 \upharpoonright #2}
\providecommand{\sto}{\rightarrow}
\providecommand{\FmlaSet}[1]{\mathrm{Fmla}_{#1}}
\providecommand{\Thm}{\mathrm{Thm}}
\providecommand{\Eval}{\operatorname{val}}
\providecommand{\blank}{\sqcup}
\providecommand{\blankTape}{\sqcup}
\providecommand{\Consis}{\operatorname{Con}}
\providecommand{\drs}{\mathrm{d}}
\providecommand{\VRule}{\mid}
\providecommand{\Bet}{\mathrm{B}}
\providecommand{\Seg}[1]{\mathrm{Seg}_{#1}}
\providecommand{\Enums}{\mathrm{En}}
\providecommand{\Z}{\mathbb{Z}}
\providecommand{\N}{\mathbb{N}}
\providecommand{\abs}[1]{\lvert #1 \rvert}
\providecommand{\floor}[1]{\lfloor #1 \rfloor}
\providecommand{\ceil}[1]{\lceil #1 \rceil}

% ===================== OLP compatibility layer (4) =====================
% --- fourth compatibility block: harvested from the per-candidate compile logs ---
\providecommand{\beredpar}{\mathrel{\rightarrow_\beta}}
\makeatletter
\def\bered{\@ifnextchar[{\bered@i}{\mathrel{\rightarrow_\beta}}}
\def\bered@i[#1]{\mathrel{\rightarrow_\beta}}
\def\bred{\@ifnextchar[{\bred@i}{\mathrel{\rightarrow_\beta}}}
\def\bred@i[#1]{\mathrel{\rightarrow_\beta}}
\makeatother
\providecommand{\beredone}{\mathrel{=_\beta}}
\providecommand{\bredone}{\mathrel{=_\beta}}
\providecommand{\xredone}{\mathrel{=_\beta}}
\providecommand{\aconvone}{\mathrel{=_\alpha}}
\providecommand{\cardless}{<}
\providecommand{\tagrefs}[1]{#1}
\providecommand{\ordplus}{\mathbin{+}}
\providecommand{\funrestrictionto}[2]{#1 \upharpoonright #2}
\providecommand{\fundefined}{\mathord{\uparrow}}
\providecommand{\Trm}{\mathrm{Trm}}
\providecommand{\OProv}{\operatorname{Prov}}
\providecommand{\OPrf}{\operatorname{Prf}}
\providecommand{\aeq}{\mathrel{=}}
\providecommand{\mClass}[1]{[#1]}
\providecommand{\equivclass}[1]{[#1]}
\providecommand{\pheight}{\operatorname{ht}}
\providecommand{\phy}{\operatorname{ht}}
\providecommand{\Cut}{\mathrm{Cut}}
\providecommand{\QuantRank}{\operatorname{qr}}
\providecommand{\elemequiv}{\mathrel{\equiv}}
\providecommand{\PIso}{\mathrm{PIso}}
\providecommand{\Theory}{\mathrm{Th}}
\providecommand{\supstrict}{\subsetneq}
\providecommand{\trcl}{\operatorname{trcl}}
\providecommand{\cref}[1]{\ref{#1}}
\providecommand{\Discharge}{\mathrm{Dis}}
\providecommand{\Weakening}{\mathrm{Wk}}
\providecommand{\ZFminus}{\mathrm{ZF}^{-}}
\providecommand{\RK}{\mathrm{RK}}
\makeatletter
\renewcommand{\Dual}{\@ifnextchar[{\Dual@i}{\ifmmode\mathrm{d}\else$\mathrm{d}$\fi}}
\def\Dual@i[#1]{\ifmmode\mathrm{d}[#1]\else$\mathrm{d}[#1]$\fi}
\makeatother
\providecommand{\RightR}{\mathrm{R}}
\providecommand{\LeftR}{\mathrm{L}}
\providecommand{\Intro}{\mathrm{I}}
\providecommand{\Axiom}{\mathrm{Ax}}
\providecommand{\formula}[1]{#1}
\providecommand{\stared}{\ast}
\providecommand{\oset}[2]{\{\, #1 \mid #2 \,\}}
\providecommand{\setof}[2]{\{\, #1 \mid #2 \,\}}
\providecommand{\pow}[1]{\mathcal{P}(#1)}
\providecommand{\ordpair}[2]{\langle #1, #2 \rangle}
\providecommand{\ordtriple}[3]{\langle #1, #2, #3 \rangle}
\providecommand{\union}{\cup}
\providecommand{\inter}{\cap}
\providecommand{\seq}[1]{\langle #1 \rangle}
\providecommand{\SuchThat}{\mid}
\providecommand{\where}{\mid}
\providecommand{\TVar}{\mathrm{TVar}}
\providecommand{\Bool}{\mathrm{Bool}}
\providecommand{\Val}{\operatorname{Val}}
\providecommand{\FV}{\operatorname{FV}}
\providecommand{\TL}[1]{\mathrm{TL}}
\providecommand{\xstared}{\ast}

% ===================== OLP compatibility layer (5) =====================
% --- fifth compatibility block: rescues the quarantined OLP candidates ---
\providecommand{\ident}{\operatorname{id}}
\providecommand{\setrank}{\operatorname{rank}}
\providecommand{\cardfont}{\mathfrak{c}}
\providecommand{\Complement}{\overline}
\providecommand{\QR}{\mathrm{QR}}
\providecommand{\becd}{\ensuremath{\ulcorner}}
\providecommand{\disjointsum}{\mathbin{\oplus}}
\providecommand{\cardeq}[2]{#1 = #2}
\providecommand{\xred}{\mathrel{\rightarrow}}
\providecommand{\cardless}{<}
\providecommand{\cardplus}{\mathbin{+}}
\providecommand{\cardtimes}{\mathbin{\cdot}}
\providecommand{\cardexpo}{\mathbin{\hat{\ }}}
\providecommand{\cardneq}{\neq}
\providecommand{\cardle}{\leq}
\providecommand{\cardnless}{\not<}
% --- sixth compatibility block: names reported by real compile logs ---
\providecommand{\ORefut}{\operatorname{Refut}}
\providecommand{\ORProv}{\operatorname{Prov}}
\providecommand{\cutrank}{\operatorname{cr}}
\providecommand{\Refut}{\operatorname{Refut}}
\providecommand{\DischargeRule}{\mathrm{Dis}}
\providecommand{\Contraction}{\mathrm{C}}
\providecommand{\bottomAlignProof}{\relax}
\providecommand{\FalseCl}{\bot}
\providecommand{\red}{\mathrel{\rightarrow}}
\providecommand{\varAssign}{\sigma}
\providecommand{\defiff}{\mathrel{\leftrightarrow}}
\providecommand{\FalseInt}{\bot}
\providecommand{\fn}{\operatorname{fn}}
\providecommand{\stared}{\ast}
\providecommand{\dotted}{\cdot}
\providecommand{\Gamma}{\Gamma}
% --- sixth compatibility block: names reported by real compile logs ---
\providecommand{\ORefut}{\operatorname{Refut}}
\providecommand{\ORProv}{\operatorname{Prov}}
\providecommand{\cutrank}{\operatorname{cr}}
\providecommand{\Refut}{\operatorname{Refut}}
\providecommand{\DischargeRule}{\mathrm{Dis}}
\providecommand{\Contraction}{\mathrm{C}}
\providecommand{\bottomAlignProof}{\relax}
\providecommand{\FalseCl}{\bot}
\providecommand{\red}{\mathrel{\rightarrow}}
\providecommand{\varAssign}{\sigma}
\providecommand{\defiff}{\mathrel{\leftrightarrow}}
\providecommand{\FalseInt}{\bot}
\providecommand{\fn}{\operatorname{fn}}
\providecommand{\stared}{\ast}
\providecommand{\dotted}{\cdot}
\providecommand{\Gamma}{\Gamma}
\begin{document}

% ==== OLP file: content/lambda-calculus/syntax/abbreviated-syntax.tex ====
% Part: lambda-calculus
% Chapter: syntax
% Section: abbreviated-syntax



\olfileid{lam}{syn}{abb}
\olsection{Abbreviated Syntax}

Terms as defined in \olref[trm]{defn:term} are sometimes cumbersome to
write, so it is useful to introduce a more concise syntax. We must of
course be careful to make sure that the terms in the concise notation
also are uniquely readable.  One widely used version called
\emph{abbreviated terms} is as follows.

\begin{enumerate}
\item When parentheses are left out, application takes place from left
  to right. For example, if $M$, $N$, $P$, and~$Q$ are terms, then
  $MNPQ$ abbreviates $(((MN)P)Q)$.
\item Again, when parentheses are left out, lambda abstraction is
  given the widest scope possible. From example, $\lambd[x][MNP]$ is
  read as $(\lambd[x][MNP])$.
\item A lambda can be used to abstract multiple variables. For
  example, $\lambd[xyz][M]$ is short for
  $\lambd[x][\lambd[y][\lambd[z][M]]]$.
\end{enumerate}

For example,
\[
\lambd[xy][xxyx \lambd[z][xz]]
\]
abbreviates
\[
(\lambd[x][(\lambd[y][((((xx)y)x)(\lambd[z][(xz)]))])]).
\]

\begin{prob}
Expand the abbreviated term $\lambd[g][(\lambd[x][g (x x)])
  \lambd[x][g (x x)]]$.
\end{prob}


% ==== OLP file: content/lambda-calculus/syntax/alpha.tex ====
% Part: lambda-calculus
% Chapter: syntax
% Section: alpha



\olfileid{lam}{syn}{alp}
\olsection{$\alpha$-Conversion}

What is the relation between $\lambd[x][x]$ and $\lambd[y][y]$? They
both represent the identity function. They are, of course,
syntactically different terms. They differ only in the name of the
bound variable, and one is the result of ``renaming'' the bound
variable in the other. This is called \emph{$\alpha$-conversion}.

\begin{defn}[Change of bound variable, $\aconvone$]
  If a term $M$ contains an occurrence of $\lambd[x][N]$, $y \notin
  \FV{N}$, and $\Subst{N}{y}{x}$ is defined, then replacing this occurrence
  by 
  \begin{equation*}
    \lambd[y][\Subst{N}{y}{x}]
  \end{equation*}
  resulting in $M'$ is called a \emph{change of bound variable}, written
  as $M \redone[\alpha] M'$.
\end{defn}

\begin{defn}[Compatibility of relation]
  A relation $R$ on terms is said to be \emph{compatible}
  if it satisfies following conditions:
  \begin{enumerate}
  \item If $R N N'$ then $R \lambd[x][N] \lambd[x][N']$
  \item If $R P P'$ then $R (PQ) (P'Q)$
  \item If $R Q Q'$ then $R (PQ) (PQ')$
  \end{enumerate}
\end{defn}

Thus let's rephrase the definition:
\begin{defn}[Change of bound variable, $\aconvone$]
  \emph{Change of bound variable} ($\redone[\alpha]$) is 
  the smallest compatible relation on terms satisfying following
  condition:
  \begin{align*}
    &\lambd[x][N] \redone[\alpha] \lambd[y][\Subst{N}{y}{x}] && \text{if
      $x \neq y$, $y \notin \FV{N}$} \\
    & &&\text{and $\Subst{N}{y}{x}$ is defined}
  \end{align*}
\end{defn}

``Smallest'' here means the relation contains only pairs that are required by
compatibility and the additional condition, and nothing else. Thus
this relation can also be defined as follows:

\begin{defn}[Change of bound variable, $\aconvone$] \ollabel{defn:aconvone}
  \emph{Change of bound variable} ($\aconvone$) is inductively
  defined as follows:
  \begin{enumerate}
  \item If $N \aconvone N'$ then $\lambd[x][N] \aconvone
    \lambd[x][N']$ \ollabel{defn:aconvone1}
  \item If $P \aconvone P'$ then $(PQ) \aconvone (P'Q)$ \ollabel{defn:aconvone2}
  \item If $Q \aconvone Q'$ then $(PQ) \aconvone (PQ')$ \ollabel{defn:aconvone3}
  \item If $x \neq y$, $y \notin \FV{N}$ and $\Subst{N}{y}{x}$ is defined, then
    $\lambd[x][N] \redone[\alpha] \lambd[y][\Subst{N}{y}{x}]$.
    \ollabel{defn:aconvone4}
  \end{enumerate}
\end{defn}

The definitions are equivalent, but we leave the proof as an
exercise. From now on we will use the inductive definition.

\begin{defn}[$\alpha$-conversion, $\aconv$]
  \emph{$\alpha$-conversion} ($\aconv$) is the smallest reflexitive
  and transitive relation on terms containing $\aconvone$.
\end{defn}

As above, ``smallest'' means the relation only contains pairs required
by transitivity, and $\aconvone$, which leads to the following
equivalent definition:

\begin{defn}[$\alpha$-conversion, $\aconv$] \ollabel{defn:aconv}
  \emph{$\alpha$-conversion} ($\aconv$) is inductively defined as follows:
  \begin{enumerate}
  \item If $P \aconv Q$ and $Q \aconv R$, then $P \aconv R$.
    \ollabel{defn:aconv1}
  \item If $P \aconvone Q$, then $P \aconv Q$. \ollabel{defn:aconv2}
  \item $P \aconv P$. \ollabel{defn:aconv3}
  \end{enumerate}
\end{defn}

\begin{ex}
  $\lambd[x][f x]$ $\alpha$-converts to $\lambd[y][f
  y]$, and conversely. Informally
  speaking, they are both functions that accept an argument and
  return $f$ of that argument, refering to the environment variable~$f$.

  $\lambd[x][f x]$ does not $\alpha$-convert to $\lambd[x][g
    x]$. Informally speaking, they refer to the environment variables
  $f$ and~$g$ respectively, and this makes them different functions: they
  behave differently in environments where $f$ and $g$ are different.
\end{ex}

\begin{prob}
  Are the following pairs of terms $\alpha$-convertible?
  \begin{enumerate}
  \item $\lambd[x][\lambd[y][x]]$ and $\lambd[y][\lambd[x][y]]$
  \item $\lambd[x][\lambd[y][x]]$ and $\lambd[c][\lambd[b][a]]$
  \item $\lambd[x][\lambd[y][x]]$ and $\lambd[c][\lambd[b][a]]$
  \end{enumerate}
\end{prob}

\begin{lem}\ollabel{lem:fv-one}
  If $P \aconvone Q$ then $\FV{P} = \FV{Q}$.
\end{lem}

\begin{proof}
  By induction on the !!{derivation} of $P \aconvone Q$.
  \begin{enumerate}
  \item If the last rule is \olref{defn:aconvone4}, then $P$ is of the
    form $\lambd[x][N]$ and $Q$ of the form
    $\lambd[y][\Subst{N}{y}{x}]$, with $x \neq y$, $y \notin \FV{N}$
    and $\Subst{N}{y}{x}$ defined. We distinguish cases according to
    whether $x \in \FV{N}$:
    \begin{enumerate}
    \item If $x \in FV(N)$, then:
      \begin{align*}
        \FV{\lambd[y][\Subst{N}{y}{x}]} &=
        \FV{\Subst{N}{y}{x}} \setminus \{y\} \\
        & = ((\FV{N} \setminus \{x\}) \cup \{y\}) \setminus \{y\}
         && \text{ by \olref[sub]{thm:infv}} \\
        & = \FV{N} \setminus \{x\} \\
        & = \FV{\lambd[x][N]}
      \end{align*}
    \item If $x \notin FV(N)$, then:
      \begin{align*}
        \FV{\lambd[y][\Subst{N}{y}{x}]}
        & = FV{\Subst{N}{y}{x}} \setminus \{y\} \\
        & = \FV{N} \setminus \{x\}
         && \text{by \olref[sub]{thm:notinfv}} \\
        & = \FV{\lambd[x][N]}.
      \end{align*}
    \end{enumerate}
  \item The other three cases are left as exercises. 
  \end{enumerate}
\end{proof}

\begin{prob}
  Complete the proof of \olref[lam][syn][alp]{lem:fv-one}.
\end{prob}

\begin{lem}\ollabel{lem:inv}
  If $P \aconvone Q$ then $Q \aconvone P$.
\end{lem}

\begin{proof}
  Induction on the !!{derivation} of $P \aconvone Q$.
  \begin{enumerate}
  \item If the last rule is \olref{defn:aconvone4}, then $P$ is of the
    form $\lambd[x][N]$ and $Q$ of the form
    $\lambd[y][\Subst{N}{y}{x}]$, where $x \neq y$, $y \notin \FV{N}$
    and $\Subst{N}{y}{x}$ defined. First, we have $y \notin
    \FV{\Subst{N}{y}{x}}$ by \olref[sub]{thm:clr}. By
    \olref[sub]{thm:inv} we have that $\Subst{\Subst{N}{y}{x}}{x}{y}$ is
    not only defined, but also equal to~$N$. Then by
    \olref{defn:aconvone4}, we have $\lambd[y][\Subst{N}{y}{x}]
    \aconvone \lambd[x][\Subst{\Subst{N}{y}{x}}{x}{y}] =
    \lambd[x][N]$.
  \end{enumerate}
\end{proof}

\begin{prob}
  Complete the proof of \olref[lam][syn][alp]{lem:inv}
\end{prob}

\begin{thm}
  $\alpha$-Conversion is an equivalence relation on terms, i.e., it is
  reflexive, symmetric, and transitive.
\end{thm}

\begin{proof}
  \begin{enumerate}
  \item For each term $M$, $M$ can be changed to $M$ by
    \emph{zero} changes of bound variables.
  \item If $P$ is $\alpha$-converts to $Q$ by a series of changes
    of bound variables, then from $Q$ we can just inverse these
    changes (by \olref{lem:inv}) in
    opposite order to obtain~$P$.
  \item If $P$ $\alpha$-converts to $Q$ by a series of changes of
    bound variables, and $Q$ to $R$ by another series, then we can
    change $P$ to $R$ by first applying the first series and then the
    second series.
  \end{enumerate}
\end{proof}

From now on we say that $M$ and $N$ are \emph{$\alpha$-equivalent}, $M
\aeq N$, iff $M$ $\alpha$-converts to $N$ (which, as we've just shown,
is the case iff $N$ $\alpha$-converts to~$M$).

\begin{thm}\ollabel{thm:fv}
  If $M \aeq N$, then $\FV{M} = \FV{N}$.
\end{thm}

\begin{proof}
  Immediate from \olref{lem:fv-one}.
\end{proof}

\begin{lem}\ollabel{lem:sub:R}
  If $R \aeq R'$ and $\Subst{M}{R}{y}$ is defined, then $\Subst{M}{R'}{y}$ is
  defined and $\alpha$-equivalent to $\Subst{M}{R}{y}$.
\end{lem}

\begin{proof}
  Exercise.
\end{proof}

\begin{prob}
  Prove \olref[lam][syn][alp]{lem:sub:R}.
\end{prob}

Recall that in \olref[sub]{sec}, substitution is undefined in some
cases; however, using $\alpha$-conversion on terms, we can make
substitution always defined by renaming bound variables. The result
preserves $\alpha$-equivalence, as shown in this theorem:

\begin{thm}\ollabel{thm:sub}
  For any $M$, $R$, and~$y$, there exists $M'$ such that $M \aeq M'$
  and $\Subst{M'}{R}{y}$ is defined. Moreover, if there is another
  pair $M'' \aeq M$ and $R''$ where $\Subst{M''}{R''}{y}$ is defined
  and $R'' \aeq R$, then $\Subst{M'}{R}{y} \aeq \Subst{M''}{R''}{y}$.
\end{thm}

\begin{proof}
  By induction on the formation of $M$:
  \begin{enumerate}
  \item $M$ is a variable~$z$: Exercise.
  \item Suppose $M$ is of the form $\lambd[x][N]$.  Select a variable
    $z$ other than $x$ and $y$ and such that $z \notin \FV{N}$ and $z
    \notin \FV{R}$.  By inductive hypothesis, we there is $N'$ such
    that $N' \aeq N$ and $\Subst{N'}{z}{x}$ is defined. Then
    $\lambd[x][N] \aeq \lambd[x][N']$ too, by
    \olref{defn:aconvone}\olref{defn:aconvone1}.  Now $\lambd[x][N']
    \aeq \lambd[z][\Subst{N'}{z}{x}]$ by
    \olref{defn:aconvone}\olref{defn:aconvone4}.  We can do this
    because $z \ne x$, $z \notin FV(N')$ and $\Subst{N'}{z}{x}$ is
    defined. Finally, $\Subst{\lambd[z][\Subst{N'}{z}{x}]}{R}{y}$ is
    defined, because $z \neq y$ and $z \notin FV(R)$.

    Moreover, if there is another $N''$ and $R''$ satisfying the same
    conditions,
    \begin{multline*}
      \Subst{(\lambd[z][\Subst{N''}{z}{x}])}{R''}{y} =\\
    \begin{aligned}
      &= \lambd[z][\Subst{\Subst{N''}{z}{x}}{R''}{y}] \\
      &= \lambd[z][\Subst{\Subst{N''}{z}{x}}{R}{y}] && \text{by
        \olref{lem:sub:R}}\\
      &=\lambd[z][\Subst{\Subst{N'}{z}{x}}{R}{y}]
      && \text{by inductive
        hypothesis}\\
      &=\Subst{(\lambd[z][\Subst{N'}{z}{x}])}{R}{y}
    \end{aligned}
    \end{multline*}
    \item $M$ is of the form $(PQ)$: Exercise.
  \end{enumerate}
\end{proof}

\begin{prob}
  Complete the proof of \olref[lam][syn][alp]{thm:sub}.
\end{prob}

\begin{cor}\ollabel{cor:sub}
  For any $M$, $R$, and~$y$, there exists a pair of $M'$ and $R'$
  such that $M' \aeq M$, $R \aeq R'$ and $\Subst{M'}{R'}{y}$ is
  defined. Moreover, if there is another pair $M'' \aeq M$ and $R''$
  with $\Subst{M'}{R'}{y}$ defined, then $\Subst{M'}{R'}{y} \aeq
  \Subst{M''}{R''}{y}$.
\end{cor}

\begin{proof}
  Immediate from \olref{thm:sub}.
\end{proof}



% ==== OLP file: content/lambda-calculus/syntax/beta.tex ====
% Part: lambda-calculus
% Chapter: introduction
% Section: beta



\olfileid{lam}{int}{bet}
\olsection{$\beta$-reduction}

When we see $(\lambd[m][(\lambd[y][y]) m])$, it is natural to
conjecture that it has some connection with $\lambd[m][m]$, namely the
second term should be the result of ``simplifying'' the first.  The
notion of \emph{$\beta$-reduction} captures this intuition formally.

\begin{defn}[$\beta$-contraction, $\bredone$] \ollabel{defn:betacontr}
  \emph{The $\beta$-contraction} ($\bredone$) is the smallest compatible
  relation on terms satisfying the following condition:
  \[
    (\lambd[x][N])Q \bredone \Subst{N}{Q}{x} 
  \]
  We say $P$ is \emph{$\beta$-contracted} to $Q$ if $P \bredone Q$.  A
  term of the form $(\lambd[x][N])Q$ is called a \emph{redex}.
\end{defn}

\begin{prob} \ollabel{prob:def}
  Spell out the equivalent inductive definitions of $\beta$-contraction as we
  did for change of bound variable in \olref[lam][syn][alp]{defn:aconvone}.
\end{prob}
  
\begin{defn}[$\beta$-reduction, $\bred$] \ollabel{defn:betared}
  \emph{$\beta$-reduction} ($\bred$) is the smallest reflexive,
  transitive relation on terms containing $\bredone$.  We say $P$ is
  $\beta$-reduced to $Q$ if $P \bred Q$.
\end{defn}

We will write $\redone$ instead of $\bredone$, and $\red$ instead of
$\bred$ when context is clear.

Informally speaking, $M \bred N$  if and only if $M$ can be changed to
$N$ by zero or several steps of $\beta$-contraction.

\begin{defn}[$\beta$-normal]
A term that cannot be $\beta$-contracted any further is said to be
\emph{$\beta$-normal}. 
\end{defn}

If $M \bred N$ and $N$ is $\beta$-normal, then we say $N$ is a
\emph{normal form} of~$M$. One may ask if the normal form of a term is
unique, and the answer is yes, as we will see later.

Let us consider some examples.
\begin{enumerate}
\item We have
\begin{align*}
(\lambd[x][xxy]) \lambd[z][z] & \redone (\lambd[z][z])(\lambd[z][z]) y \\
& \redone (\lambd[z][z]) y \\
& \redone y
\end{align*}
\item ``Simplifying'' a term can actually make it more complex:
\begin{align*}
(\lambd[x][xxy])(\lambd[x][xxy]) & \redone (\lambd[x][xxy])(\lambd[x][xxy])y \\
& \redone (\lambd[x][xxy])(\lambd[x][xxy])yy \\
& \redone \dots
\end{align*}
\item It can also leave a term unchanged:
\[
(\lambd[x][xx])(\lambd[x][xx]) \redone (\lambd[x][xx])(\lambd[x][xx])
\]
\item Also, some terms can be reduced in more than one way; for
  example,
\[
(\lambd[x][(\lambd[y][yx]) z]) v \redone (\lambd[y][yv]) z
\]
by contracting the outermost application; and
\[
(\lambd[x][(\lambd[y][yx]) z]) v \redone (\lambd[x][zx]) v
\]
by contracting the innermost one. Note, in this case, however, that
both terms further reduce to the same term, $zv$.
\end{enumerate}

The final outcome in the last example is not a coincidence, but rather
illustrates a deep and important property of the lambda calculus, known as the
Church--Rosser property.

\begin{digress}
  In general, there is more than one way to $\beta$-reduce a term,
  thus many reduction strategies have been invented, among which the
  most common is the \emph{natural strategy}.  The natural strategy
  always contracts the \emph{left-most} redex, where the position of a
  redex is defined as its starting point in the term.  The natural
  strategy has the useful property that a term can be reduced to a
  normal form by some strategy iff it can be reduced to normal form
  using the natural strategy. In what follows we will use the natural
  stratuegy unless otherwise specified.
\end{digress}

\begin{defn}[$\beta$-equivalence, $\equal$]
  \emph{$\beta$-Equivalence} ($\equal$) is the relation inductively
  defined as follows:
  \begin{enumerate}
  \item $M \equal M$.
  \item If $M \equal N$, then $N \equal M$.
  \item If $M \equal N$, $N \equal O$, then $M \equal O$.
  \item If $M \equal N$, then $PM \equal PN$.
  \item If $M \equal N$, then $MQ \equal NQ$.
  \item If $M \equal N$, then $\lambd[x][M] \equal \lambd[x][N]$.
  \item $(\lambd[x][N])Q \equal \Subst{N}{Q}{x}$.
  \end{enumerate}
\end{defn}

The first three rules make the relation an equivalence relation; the
next three make it compatible; the last ensures that it contains
$\beta$-contraction.

Informally speaking, two terms are $\beta$-equivalent if and only if
one of them can be changed to the other in zero or more steps of
$\beta$-contraction, or ``inverse'' of $\beta$-contraction. The inverse of
$\beta$-contraction is defined so that $M$ inverse-$\beta$-contracts to $N$
iff $N$ $\beta$-contracts to $M$.

Besides the above rules, we will extend the relation with more
rules, and denote the extended equivalence relation as $\equal[X]$,
where $X$ is the extending rule.


% ==== OLP file: content/lambda-calculus/syntax/conversion.tex ====
% Part: lambda-calculus
% Chapter: introduction
% Section: conversion



\olfileid{lam}{int}{con}
\olsection{Conversion and Reduction}

There three type of conversions or reductions in lambda calculus. The
difference between conversion and reduction is whether the
``transformation'' is unidirectional or bidirectional, as we will see.

The $\alpha$-conversion is about variable renaming. It is natural to
consider $\lambda x. x$ and $\lambda y.y$ representing the same
function (the identity function), we formallize this idea as follows:




% ==== OLP file: content/lambda-calculus/syntax/de-bruijn.tex ====
% Part: lambda-calculus
% Chapter: syntax
% Section: de-bruijn



\olfileid{lam}{syn}{deb}

\olsection{The De Bruijn Index}

$\alpha$-Equivalence is very natural, as terms that are
$\alpha$-equivalent ``mean the same.'' In fact, it is possible to give
a syntax for lambda terms which does not distinguish terms that can be
$\alpha$-converted to each other. The best known replaces variables by
their \emph{De Bruijn index}.

When we write $\lambd[x][M]$, we explicitly state that $x$ is the
parameter of the function, so that we can use $x$ in $M$ to refer to
this parameter. In the de Bruijn index, however, parameters have no
name and reference to them in the function body is denoted by a number
denoting the levels of abstraction between them. For example, consider
the example of $\lambd[x][\lambd[y][y x]]$: the outer abstraction is on
binds the variable~$x$; the inner abstraction binds the variable
is~$y$; the sub-term $y x$ lies in the scope of the inner abstraction:
there is no abstraction between $y$ and its abstract~$\lambd[y]$, but
one abstract between $x$ and its abstract~$\lambd[x]$. Thus we write
$0\, 1$ for $y x$, and $\lambd[][\lambd[][01]]$ for the entire term.

\begin{defn}
  De Bruijn terms are inductively defines as follows:
  \begin{enumerate}
  \item $n$, where $n$ is any natural number.
  \item $PQ$, where $P$ and $Q$ are both De Bruijn terms.
  \item $\lambd[][N]$, where $N$ is a De Bruijn term.
  \end{enumerate}
\end{defn}

A formalized translation from ordinary lambda terms to De Bruijn
indexed terms is as follows:
\begin{defn}
  \begin{align*}
    F_\Gamma(x) &= \Gamma(x) \\
    F_\Gamma(PQ) &= F_\Gamma(P)F_\Gamma(Q) \\
    F_\Gamma(\lambd[x][N]) &= \lambd[][F_{x,\Gamma}(N)]
  \end{align*}
  where $\Gamma$ is a list of variables indexed from zero, and
  $\Gamma(x)$ denotes the position of the variable $x$ in $\Gamma$.
  For example, if $\Gamma$ is $x,y,z$, then $\Gamma(x)$ is $0$ and
  $\Gamma(z)$ is $2$.
  
  $x,\Gamma$ denotes the list resulted from pushing $x$ to the head of
  $\Gamma$; for instance, continuing the $\Gamma$ in last example,
  $w,\Gamma$ is $w,x,y,z$.
\end{defn}

Recovering a standard lambda term from a de Bruijn term is done as
follows:

\begin{defn}
  \begin{align*}
    G_\Gamma(n) &= \Gamma[n] \\
    G_\Gamma(PQ) &= G_\Gamma(P) G_\Gamma(Q) \\
    G_\Gamma(\lambd[][N]) &= \lambd[x][G_{x,\Gamma}(N)]
  \end{align*}
  where $\Gamma$ is again a list of variables indexed from zero, and
  $\Gamma[n]$ denotes the variable in position~$n$. For example,
  if $\Gamma$ is $x,y,z$, then $\Gamma[1]$ is $y$.

  The variable $x$ in last equation is chosen to be any variable that
  not in $\Gamma$.
\end{defn}

Here we give some results without proving them:

\begin{prop}
  If $M \aconvone M'$, and $\Gamma$ is any list containing $\FV{M}$, then
  $F_\Gamma(M) \eqs F_\Gamma(M')$.
\end{prop}


% ==== OLP file: content/lambda-calculus/syntax/eta.tex ====
% Part: lambda-calculus
% Chapter: syntax
% Section: eta



\olfileid{lam}{syn}{eta}
\olsection{$\eta$-conversion}

There is another relation on $\lambda$ terms. In \olref[fv]{sec} we
used the example $\lambd[x][(fx)]$, which accepts an argument and
applies $f$ to it.  In other words, it is the same function as~$f$:
$\lambd[x][(fx)]N$ and $fN$ both reduce to~$fN$.  We use
$\eta$-reduction (and $\eta$-extension) to capture this idea.

\begin{defn}[$\eta$-contraction, $\eredone$]
  \ollabel{defn:beredone}
  \emph{$\eta$-contraction} ($\eredone$) is the smallest compatible relation
  on terms satisfying the following condition:
  \[
    \lambd[x][M x] \eredone M \text{ provided } x \notin FV(M)
  \]
\end{defn}

\begin{defn}[$\beta\eta$-reduction, $\bered$] \ollabel{defn:bered}
  \emph{$\beta\eta$-reduction} ($\bered$) is the smallest reflexive,
  transitive relation on terms containing $\bredone$ and $\eredone$,
  i.e., the rules of reflexivity and transitive plus the following two
  rules:
  \begin{enumerate}
  \item If $M \bredone N$ then $M \bered N$. \ollabel{defn:bered3}
  \item If $M \eredone N$ then $M \bered N$. \ollabel{defn:bered4}
  \end{enumerate}
    
\end{defn}

\begin{defn}
  We extend the equivalence relation $\equal$ with the $\eta$-conversion rule:
  \[
  \lambd[x][f x] \equal f
  \]
  and denote the extended relation as $\equal[\eta]$.
\end{defn}

$\eta$-equivalence is important because it is related to extensionality
of lambda terms:

\begin{defn}[Extensionality]
  We extend the equivalence relation $\equal$ with the (\ext) rule:
  \begin{center}
  If $Mx \equal Nx$ then $M \equal N$, provided $x \notin FV(MN)$.
  \end{center}
  and denote the extended relation as $\equal[\ext]$.
\end{defn}

Roughly speaking, the rule states that two terms, viewed as functions,
should be considered equal if they behave the same for the same
argument.

We now prove that the $\eta$ rule provides exactly the extensionality,
and nothing else.

\begin{thm}
  $M \equal[\ext] N$ if and only if $M \equal[\eta] N$.
\end{thm}

\begin{proof}
  First we prove that $\equal[\eta]$ is closed under the
  extensionality rule. That is, $ext$ rule doesn't add anything to
  $\equal[\eta]$. We then have $\equal[\eta]$ contains $\equal[\ext]$,
  and if $M \equal[\ext] N$, then $M \equal[\eta] N$.

  To prove $\equal[\eta]$ is closed under \ext, note that for any $M
  \equal N$ derived by the \ext{} rule, we have $Mx
  \equal[\eta] Nx$ as premise. Then we have $\lambd[x][Mx]
  \equal[\eta] \lambd[x][Nx]$ by a rule of $\equal$, applying $\eta$
  on both side gives us $M \equal[\eta] N$.

  Similarly we prove that the $\eta$ rule is contained in
  $\equal[ext]$. For any $\lambd[x][Mx]$ and $M$ with $x \notin
  FV(M)$, we have that $(\lambd[x][Mx])x \equal[\ext] Mx$,
  giving us $\lambd[x][Mx] \equal[\ext] M$ by the \ext{} rule.
\end{proof}


% ==== OLP file: content/lambda-calculus/syntax/free-variables.tex ====
% Part: lambda-calculus
% Chapter: introduction
% Section: free-variables



\olfileid{lam}{syn}{fv}
\olsection{Free Variables}

Lambda calculus is about functions, and lambda abstraction is how
functions arise. Intuitively, $\lambd[x][M]$ is the function with
values given by~$M$ when the argument to the function is assigned
to~$x$. But not every occurrence of~$x$ in~$M$ is relevant: if $M$
contains another abstract $\lambd[x][N]$ then the occurrences of $x$
in~$N$ are relevant to $\lambd[x][N]$ but not to $\lambd[x][M]$.  So,
a lambda abstract $\lambd[x]$ inside $\lambd[x][M]$ \emph{binds} those
occurrences of $x$ in~$M$ that are not already bound by another lambda
abstract---the \emph{free} occurrences of~$x$ in~$M$.

\begin{defn}[Scope]
If $\lambd[x][M]$ occurs inside a term~$N$, then the corresponding
occurrence of~$N$ is the \emph{scope} of the~$\lambd[x]$.
\end{defn}


\begin{defn}[Free and bound occurrence]
  An occurrence of variable $x$ in a term~$M$ is \emph{free} if it is
  not in the scope of a $\lambd[x]$, and \emph{bound} otherwise. An
  occurrence of a variable~$x$ in $\lambd[x][M]$ is bound by the initial
  $\lambd[x]$ iff the occurrence of $x$ in~$M$ is free.
\end{defn}

\begin{ex}
  In $\lambd[x][x y]$, both $x$ and $y$ are in the scope of $\lambd[x]$,
  so $x$ is bound by $\lambd[x]$. Since $y$ is not in the scope of
  any $\lambd[y]$, it is free.  In $\lambd[x][x x]$, both occurrences
  of $x$ are bound by~$\lambd[x]$, since both are free in $xx$.  In
  $((\lambd[x][xx])x)$, the last occurrence of~$x$ is free, since it
  is not in the scope of a $\lambd[x]$. In
  $\lambd[x][(\lambd[x][x])x]$, the scope of the first $\lambd[x]$ is
  $(\lambd[x][x])x$ and the scope of the second $\lambd[x]$ is the
  second-to-last occurrence of~$x$. In $(\lambd[x][x])x$, the last
  occurrence of $x$ is free, and the second-to-last is bound.  Thus,
  the second-to-last occurrence of $x$ in $\lambd[x][(\lambd[x][x])x]$
  is bound by the second $\lambd[x]$, and the last occurrence by the
  first $\lambd[x]$.
\end{ex}

For a term $P$, we can check all variable occurrences in it and get a
set of free variables. This set is denoted by $\FV{P}$ with a natural
definition as follows:

\begin{defn}[Free variables of a term] \ollabel{def:fv}
  The set of \emph{free variables} of a term is defined inductively by:
  \begin{enumerate}
    \item $\FV{x} = \{x\}$ \ollabel{def:fv1} 
    \item $\FV{\lambd[x][N]} = \FV{N} \setminus \{x\}$    \ollabel{def:fv2}
    \item $\FV{PQ} = \FV{P} \cup \FV{Q}$ \ollabel{def:fv3}
  \end{enumerate}
\end{defn}

\begin{prob}
  \begin{enumerate}
  \item Identify the scopes of $\lambd[g]$ and the two $\lambd[x]$ in
    this term: $\lambd[g][(\lambd[x][g (x x)]) \lambd[x][g (x x)]]$.
  \item In $\lambd[g][(\lambd[x][g (x x)]) \lambd[x][g (x x)]]$, are
    all occurrences of variables bound? By which abstractions are they
    bound respectively?
    \item Give $\FV{\lambd[x][(\lambd[y][(\lambd[z][x y]) z]) y]}$
  \end{enumerate}
\end{prob}

\begin{explain}
A free variable is like a reference to the outside world (the
\emph{environment}), and a term containing free variables can be seen
as a partially specified term, since its behaviour depends on how we
set up the environment. For example, in the term $\lambd[x][f x]$, which
accepts an argument~$x$ and returns $f$ of that argument, the
variable~$f$ is free. This value of the term is dependent on the
environment it is in, in particular the value of $f$ in that
environment.

If we apply abstraction to this term, we get $\lambd[f][\lambd[x][f
    x]]$. This term is no longer dependent on the environment variable
$f$, because it now designates a function that accepts two arguments
and returns the result of applying the first to the second. Changing
$f$ in the environment won't have any effect on the behavior of this
term, as the term will only use whatever is passed as an argument, and
not the value of $f$ in the environment.
\end{explain}

\begin{defn}[Closed term, combinator]
  A term with no free variables is called a \emph{closed term}, or a
  \emph{combinator}.
\end{defn}

\begin{lem}\ollabel{lem:fv}
  \begin{enumerate}
    \item \ollabel{lem:fv-abs} If $y \neq x$, then $y \in
      \FV{\lambd[x][N]}$ iff $y \in \FV{N}$.
    \item \ollabel{lem:fv-app} $y \in \FV{PQ}$ iff $y \in \FV{P}$ or $y
      \in \FV{Q}$.
    \end{enumerate}
\end{lem}

\begin{proof}
  Exercise.
\end{proof}

\begin{prob}
  Prove \olref[lam][syn][fv]{lem:fv}.
\end{prob}


% ==== OLP file: content/lambda-calculus/syntax/substitution.tex ====
% Part: lambda-calculus
% Chapter: syntax
% Section: substitution



\olfileid{lam}{syn}{sub}
\olsection{Substitution}

\begin{explain}
Free variables are references to environment variables, thus it makes
sense to actually use a specific value in the place of a free
variable. For example, we may want to replace $f$ in $\lambd[x][f x]$
with a specific term, like the identity function $\lambd[y][y]$. This
results in $\lambd[x][(\lambd[y][y]) x]$. The process of replacing
free variables with lambda terms is called substitution.
\end{explain}

\begin{defn}[Substitution] \ollabel{defn:substitution}
  The \emph{substitution} of a term $N$ for a variable $x$ in a term
  $M$, $\Subst{M}{N}{x}$, is defined inductively by:
  \begin{enumerate}
    \item $\Subst{x}{N}{x} = N$. \ollabel{defn:substitution-1} 
    \item $\Subst{y}{N}{x}  = y$ if $x \neq y$. \ollabel{defn:substitution-2} 
    \item $\Subst{PQ}{N}{x}  = (\Subst{P}{N}{x}) (\Subst{Q}{N}{x})$.
      \ollabel{def:substitution-3}
    \item $\Subst{(\lambd[y][P])}{N}{x} = \lambd[y][\Subst{P}{N}{x}]$,
      if $x \neq y$ and $y \notin \FV{N}$, otherwise undefined.
      \ollabel{defn:substitution-4}
  \end{enumerate}
\end{defn}

\begin{explain}
In \olref{defn:substitution}\olref{defn:substitution-4}, we require $x
\neq y$ because we don't want to replace \emph{bound} occurrences of
the variable~$x$ in~$M$ by~$N$.  For example, if we compute the
substitution $\Subst{\lambd[x][x]}{y}{x}$, the result should not be
$\lambd[x][y]$ but simply $\lambd[x][x]$.

When substituting $N$ for~$x$ in $\lambd[y][P]$, we also require that
$y \notin \FV{N}$.  For example, we cannot substitute $y$ for $x$ in
$\lambd[y][x]$, i.e., $\Subst{\lambd[y][x]}{y}{x}$, because it would
result in $\lambd[y][y]$, a term that stands for the function that
accepts an argument and returns it directly. But the term
$\lambd[y][x]$ stands for a function that always returns the term~$x$
(or whatever $x$ refers to). So the result we actually want is a
function that accepts an argument, drop it, and returns the
environment variable~$y$. To do this properly, we would first have to
``rename'' the bound variable~$y$.
\end{explain}

\begin{prob}
  What is the result of the following substitutions?
  \begin{enumerate}
  \item $\Subst{\lambd[y][x(\lambd[w][vwx])]}{(uv)}{x}$
  \item $\Subst{\lambd[y][x(\lambd[x][x])]}{(\lambd[y][xy])}{x}$
  \item $\Subst{y(\lambd[v][xv])}{(\lambd[y][vy])}{x}$
  \end{enumerate}
\end{prob}

\begin{thm} \ollabel{thm:notinfv}
  If $x \notin \FV{M}$, then $\FV{\Subst{M}{N}{x}} = \FV{M}$, if
  the left-hand side is defined.
\end{thm}

\begin{proof}
  By induction on the formation of $M$.
  \begin{enumerate}
  \item $M$ is a variable: exercise.
  \item $M$ is of the form $(PQ)$: exercise.
  \item $M$ is of the form $\lambd[y][P]$, and since
    $\Subst{\lambd[y][P]}{N}{x}$ is defined, it has to be
    $\lambd[y][\Subst{P}{N}{x}]$. Then $\Subst{P}{N}{x}$ has to be
    defined; also, $x \neq y$ and $x \notin \FV{Q}$. Then:
    \begin{multline*}
      \FV{\Subst{\lambd[y][P]}{N}{x}} = \\
    \begin{aligned}
      & = \FV{\lambd[y][\Subst{P}{N}{x}]} &&
      \text{by \olref{defn:substitution-4}}\\
      & = \FV{\Subst{P}{N}{x}} \setminus \{y\} &&
      \text{by \olref[fv]{def:fv}\olref[fv]{def:fv2}}\\
      & = \FV{P} \setminus \{y\} && \text{by inductive hypothesis} \\
      & = \FV{\lambd[y][P]} && \text{by \olref[fv]{def:fv}\olref[fv]{def:fv2}}
    \end{aligned}
    \end{multline*}
  \end{enumerate}
\end{proof}

\begin{prob}
  Complete the proof of \olref[lam][syn][sub]{thm:notinfv}.
\end{prob}

\begin{thm} \ollabel{thm:infv}
  If $x \in \FV{M})$, then $\FV{\Subst{M}{N}{x}} = (\FV{M} \setminus
  \{x\}) \cup \FV{N}$, provided the left hand is defined.
\end{thm}

\begin{proof}
  By induction on the formation of $M$.
  \begin{enumerate}
  \item $M$ is a variable: exercise.
  \item $M$ is of the form $PQ$: Since
    $\Subst{(PQ)}{N}{y}$ is defined, it has to be
    $(\Subst{P}{N}{x})(\Subst{Q}{N}{x})$ with both substitution
    defined. Also, since $x \in \FV{PQ}$, either $x \in \FV{P}$ or
    $x \in \FV{Q}$ or both. The rest is left as an exercise.
  \item $M$ is of the form $\lambd[y][P]$. Since
    $\Subst{\lambd[y][P]}{N}{x}$ is defined, it has to be
    $\lambd[y][\Subst{P}{N}{x}]$, with $\Subst{P}{N}{x}$
    defined, $x \neq y$ and $y \notin \FV{N}$; also, since $y \in
    \FV{\lambd[x][P]}$, we have $y \in \FV{P}$ too. Now:
    \begin{multline*}
      \FV{\Subst{(\lambd[y][P])}{N}{x}} = \\
      \begin{aligned}
      & = \FV{\lambd[y][\Subst{P}{N}{x}]} \\
      & = \FV{\Subst{P}{N}{x}} \setminus \{y\} \\
      & = ((\FV{P} \setminus \{y\}) \cup (\FV{N} \setminus \{x\})
       && \text{by inductive hypothesis} \\
      & = (\FV{P} \setminus \{x, y\}) \cup \FV{N}
       && x \notin \FV{N} \\
      & = (\FV{\lambd[y][P]} \setminus \{x\}) \cup \FV{N}
      \end{aligned}
      \end{multline*}
  \end{enumerate}
\end{proof}

\begin{prob}
  Complete the proof of \olref[lam][syn][sub]{thm:infv}.
\end{prob}

\begin{thm}\ollabel{thm:clr}
  $x \notin \FV{\Subst{M}{N}{x}}$, if the right-hand side is
  defined and $x \notin \FV{N}$.
\end{thm}

\begin{proof}
  Exercise.
\end{proof}

\begin{prob}
  Prove \olref[lam][syn][sub]{thm:clr}.
\end{prob}


\begin{thm}\ollabel{thm:inv}
  If $\Subst{M}{y}{x}$ is defined and $y \notin \FV{M}$, then
  $\Subst{\Subst{M}{y}{x}}{x}{y} = M$.
\end{thm}

\begin{proof}
  By induction on the formation of $M$.
  \begin{enumerate}
  \item $M$ is a variable $z$: Exercise.
  \item $M$ is of the form $(PQ)$. Then:
    \begin{align*}
      \Subst{\Subst{(PQ)}{y}{x}}{x}{y}
      &=\Subst{((\Subst{P}{y}{x})(\Subst{Q}{y}{x}))}{x}{y} \\
      &= (\Subst{\Subst{P}{y}{x}}{x}{y})(\Subst{\Subst{Q}{y}{x}}{x}{y}) \\
      &= (PQ) \text{ by inductive hypothesis}
    \end{align*}
  \item $M$ is of the form $\lambd[z][N]$. Because
    $\Subst{\lambd[z][N]}{y}{x}$ is defined, we know
    that $z \neq y$. So:
    \begin{align*}
      \Subst{\Subst{(\lambd[z][N])}{y}{x}}{x}{y}\\
      & = \Subst{(\lambd[z][\Subst{N}{y}{x}])}{x}{y} \\
      & = \lambd[z][\Subst{\Subst{N}{y}{x}}{x}{y}] \\
      &= \lambd[z][N] \text{ by inductive hypothesis} 
    \end{align*}
  \end{enumerate}
\end{proof}

\begin{prob}
  Complete the proof of \olref[lam][syn][sub]{thm:inv}.
\end{prob}


% ==== OLP file: content/lambda-calculus/syntax/syntax.tex ====
% Part: lambda-calculus
% Chapter: syntax



\olchapter{lam}{syn}{Syntax}

\olimport{terms}

\olimport{unique-readability}

\olimport{abbreviated-syntax}

\olimport{free-variables}

\olimport{substitution}

\olimport{alpha}

\olimport{de-bruijn}

\olimport{term-revisited}

\olimport{beta}

\olimport{eta}

\OLEndChapterHook


% ==== OLP file: content/lambda-calculus/syntax/term-revisited.tex ====
% Part: lambda-calculus
% Chapter: introduction
% Section: term-revisited



\olfileid{lam}{syn}{tr}
\olsection{Terms as $\alpha$-Equivalence Classes}

From now on, we will consider terms up to $\alpha$-equivalence.  That
means when we write a term, we mean its $\alpha$-equivalance class it
is in. For example, we write $\lambd[a][\lambd[b][a c]]$ for the set
of all terms $\alpha$-equivalent to it, such as $\lambd[a][\lambd[b][a
    c]]$, $\lambd[b][\lambd[a][b c]]$, etc.

Also, while in previous sections letters such as $N, Q$ are used to
denote a term, from now on we use them to denote a class, and it is
these classes instead of terms that will be our subjects of study in
what follows. Letters such as $x, y$ continues to denote a variable.

We also adopt the notation $\rep{M}$ to denote an arbitrary
!!{element} of the class~$M$, and $\rep{M}[0], \rep{M}[1], etc. $ if
we need more than one.

We reuse the notations from terms to simplify our wording. We
have following definition on classes:
\begin{defn}
  \begin{enumerate}
  \item $\lambd[x][N]$ is defined as the class containing
    $\lambd[x][\rep{N}]$.
  \item $PQ$ is defined to be the class containing $\rep{P}\rep{Q}$.
  \end{enumerate}
\end{defn}

It is not hard to see that they are well defined, because
$\alpha$-conversion is compatible.

\begin{defn} \ollabel{def:fv}
  The \emph{free variables} of an $\alpha$-equivalence class~$M$, or
  $FV(M)$, is defined to be $FV(\rep{M})$.
\end{defn}

This is well defined since $FV(\rep{M}[0]) =
FV(\rep{M}[1])$, as shown in \olref[alp]{thm:fv}.

We also reuse the notation for substition into classes:
\begin{defn} \ollabel{def:sub}
  The \emph{substitution} of $R$ for $y$ in $M$, or $\Subst{M}{R}{y}$, is
  defined to be $\Subst{\rep{M}}{\rep{R}}{y}$, for any $\rep{M}$ and
  $\rep{R}$ making the substition defined.
\end{defn}

This is also well defined as shown in \olref[alp]{cor:sub}.

Note how this definition significantly simplifies our reasoning. For
example:
\begin{align}
  \Subst{\lambd[x][x]}{y}{x} & =\ollabel{eq:1}\\
  &= \Subst{\lambd[z][z]}{y}{x} \ollabel{eq:2}\\
  &= \lambd[z][\Subst{z}{y}{x}] \\
  &= \lambd[z][z]
\end{align}

\olref{eq:1} is undefined if we still regard it as substitution on
terms; but as mentioned earlier, we now consider it a substitution on
classes, which is why \olref{eq:2} can happen: we can replace
$\lambd[x][x]$ with $\lambd[z][z]$ because they belong to the same
class.

For the same reason, from now on we will assume that the
representatives we choose always satisfy the conditions needed for
substitution. For example, when we see $\Subst{\lambd[x][N]}{R}{y}$,
we will assume the representative $\lambd[x][N]$ is chosen so that
$x \neq y$ and $x \notin FV(R)$.

Since it is a bit strange to call $\lambd[x][x]$ a ``class'', let's
call them $\Lambda$-terms (or simply ``terms'' in the rest of the
part) from now on, to distinguish them from $\lambda$-terms that we
are familiar with.

\begin{editorial}
  We cannot say goodbye to terms yet: the whole definition of
  $\Lambda$-terms is based on $\lambda$-terms, and we haven't provided
  a method to define functions on $\Lambda$-terms, which means all
  such functions have to be first defined  on $\lambda$-terms, and
  then ``projected'' to $\Lambda$-terms, as we did for substitutions.
  However we assume the reader can intuitively understand how we can
  define functions on $\Lambda$-terms.
\end{editorial}

% ==== OLP file: content/lambda-calculus/syntax/terms.tex ====
% Part: lambda-calculus
% Chapter: syntax
% Section: terms



\olfileid{lam}{syn}{trm}
\olsection{Terms}

The terms of the lambda calculus are built up inductively from an
infinite supply of variables $\Obj{v_0}$, $\Obj{v_1}$, \dots, the
symbol~``$\lambd$'', and parentheses. We will use $x$, $y$, $z$, \dots{}
to designate variables, and $M$, $N$, $P$, \dots{} to desginate terms.

\begin{defn}[Terms] \ollabel{defn:term}
The set of \emph{terms} of the lambda calculus is defined inductively by:
\begin{enumerate}
  \item \ollabel{defn:term-var} If $x$ is a variable, then $x$ is a
    term.
  \item \ollabel{defn:term-abs} If $x$ is a variable and $M$ is a
    term, then $(\lambd[x][M])$ is a term.
  \item \ollabel{defn:term-app} If both $M$ and $N$ are terms, then
    $(MN)$ is a term.
\end{enumerate}
\end{defn}

If a term $(\lambd[x][M])$ is formed according to
\olref{defn:term-abs} we say it is the result of an
\emph{abstraction}, and the $x$ in $\lambd[x]$ is called a
\emph{!!{parameter}}.  A term $(MN)$ formed according to
\olref{defn:term-app} is the result of an \emph{application}.

The terms defined above are fully parenthesized. This can get rather
cumbersome, as the term
$(\lambd[x][((\lambd[x][x])(\lambd[x][(xx)]))])$ demnostrates. We will
introduce conventions for avoiding parentheses.  However, the official
definition makes it easy to determine how a term is constructed
according to \olref{defn:term}. For example, the last step of forming
the term $(\lambd[x][((\lambd[x][x])(\lambd[x][(xx)]))])$ must be
abstraction where the !!{parameter} is~$x$. It results by abstraction
from the term $((\lambd[x][x])(\lambd[x][(xx)]))$, which is an
application of two terms. Each of these two terms is the result of an
abstraction, and so on.

\begin{prob}
Describe the formation of $(\lambd[g][(\lambd[x][(g (x x))])
  (\lambd[x][(g (x x))])])$.
\end{prob}


% ==== OLP file: content/lambda-calculus/syntax/unique-readability.tex ====
% Part: lambda-calculus
% Chapter: introduction
% Section: unique-readability



\olfileid{lam}{syn}{unq}
\olsection{Unique Readability}

We may wonder if for each term there is a unique way of forming it,
and there is.  For each lambda term there is only one way to construct
and interpret it.  In the following discussion, a \emph{formation} is
the procedure of constructing a term using the formation rules (one or
several times) of \olref[trm]{defn:term}.

\begin{lem}\ollabel{lem:term-start}
A term starts with either a variable or a parenthesis.
\end{lem}

\begin{proof}
Something counts as a term only if it is constructed according to
\olref[trm]{defn:term}. If it is the result of
\olref[trm]{defn:term-var}, it must be a variable. If it is the result
of \olref[trm]{defn:term-abs} or \olref[trm]{defn:term-app}, it starts
with a parenthesis.
\end{proof}

\begin{lem}\ollabel{lem:app-start}
The result of an application starts with either two parentheses or a
parenthesis and a variable.
\end{lem}

\begin{proof}
If $M$ is the result of an application, it is of the form $(PQ)$, so
it begins with a parenthesis. Since $P$ is a term, by
\olref{lem:term-start}, it begins either with a parenthesis or a
variable.
\end{proof}

\begin{lem}\ollabel{lem:initial}
No proper initial part of a term is itself a term.
\end{lem}

\begin{prob}
Prove \olref[lam][syn][unq]{lem:initial} by induction on the length
of terms.
\end{prob}

\begin{prop}[Unique Readability] \ollabel{prop:unq}
There is a unique formation for each term. In other words, if a term
$M$ is formed by a formation, then it is the only formation that can
form this term.
\end{prop}

\begin{proof}
  We prove this by induction on the formation of terms. 
  \begin{enumerate}
    \item $M$ is of the form
      $x$, where $x$ is some variable. Since the results of abstractions and applications always start with parentheses, they cannot have been used to
      construct $M$; Thus, the formation of $M$ must be a single step
      of \olref[trm]{defn:term}\olref[trm]{defn:term-var}.
    \item $M$ is of the form $(\lambd[x][N])$, where $x$ is some
      variable and $N$ is a term. It could not have been constructed
      according to \olref[trm]{defn:term}\olref[trm]{defn:term-var},
      because it is not a single variable.  It is not the result of an
      application, by \olref{lem:app-start}. Thus $M$ can only be the
      result of an abstraction on~$N$.  By inductive hypothesis we
      know that formation of $N$ is itself unique.
    \item $M$ is of the form $(PQ)$, where $P$ and $Q$ are terms.
      Since it starts with a parentheses, it cannot also be
      constructed by
      \olref[trm]{defn:term}\olref[trm]{defn:term-var}.  By
      \olref{lem:term-start}, $P$ cannot begin with $\lambd$, so
      $(PQ)$ cannot be the result of an abstraction.  Now suppose
      there were another way of constructing $M$ by application, e.g.,
      it is also of the form $(P'Q')$. Then $P$ is a proper initial
      segment of $P'$ (or vice versa), and this is impossible by
      \olref{lem:initial}. So $P$ and $Q$ are uniquely
      determined, and by inductive hypothesis we know that formations
      of $P$ and $Q$ is unique.
  \end{enumerate}
\end{proof}

A more readable paraphrase of the above proposition is as follows:
\begin{prop}
  A term $M$ can only be one of the following forms:
  \begin{enumerate}
    \item $x$, where $x$ is a variable uniquely determined by $M$.
    \item $(\lambd[x][N])$, where $x$ is a variable and $N$ is
      another term, both of which is uniquely determined by $M$.
    \item $(PQ)$, where $P$ and $Q$ are two terms uniquely
      determined by $M$.
  \end{enumerate}
\end{prop}


